% !TeX root = booklet-01.tex
\section{از ولتاژ فیزیکی به منطق دیجیتال}
\subsection{چرا مقدار منطقی با ولتاژ یکی نیست؟}
سیگنال فیزیکی مدار، ولتاژ یا جریان است؛ مدل دیجیتال بخشی از این رفتار را با
مقادیر گسسته صفر و یک توصیف می‌کند. سیگنال \lr{Analog} می‌تواند در یک بازه
مقادیر پیوسته داشته باشد؛ در توصیف \lr{Digital}، بازه‌هایی از ولتاژ به مقدارهای
منطقی نگاشت می‌شوند. خود ولتاژ واقعی مدار دیجیتال نیز هنگام تغییر پیوسته است.
در قرارداد این درس، ولتاژ نزدیک \lr{Ground} نماینده صفر و نزدیک
\(V_{\mathrm{DD}}\) نماینده یک است \cite[بخش \secnum{۱.۳}]{weste}.
نماد \(V_{\mathrm{DD}}\) نام ریل تغذیه مثبت است، نه یک مقدار ثابت همگانی.

\textbf{مثال آموزشی با مرزهای فرضی.} فرض کنید در مداری با تغذیه
\(\qty{1}{\volt}\)، ورودی از صفر تا \(\qty{0.3}{\volt}\) صفر معتبر و
از \(\qty{0.7}{\volt}\) تا تغذیه یک معتبر تعریف شود.
ولتاژ \(\qty{0.2}{\volt}\) صفر معتبر است، هرچند صفر ولت نیست؛
\(\qty{0.8}{\volt}\) یک معتبر و \(\qty{0.5}{\volt}\) در بازه نامعین است.
این مرزها صرفاً برای آموزش‌اند و مشخصات یک فناوری واقعی را بیان نمی‌کنند.
سطح‌های معتبر و \NotesTopic{sec:noise-margin}{\lr{Noise Margin}} از
مشخصه انتقال مدار استخراج می‌شوند.

\subsection{بیت، بردار و شمارش حالت‌ها}
یک \lr{Bit} یک مقدار دودویی دارد. دو بیت مستقل چهار ترکیب و \(n\) بیت مستقل
\(2^n\) ترکیب دارند؛ هر بیت تعداد حالت‌ها را دو برابر می‌کند.
برای سه ورودی، بنابراین هشت سطر در جدول کامل نیاز داریم.
برای آشنایی با نمایش عددی، سه بیت \(b_2,b_1,b_0\) در نمایش بدون علامت، عدد
\(4b_2+2b_1+b_0\) را نمایش می‌دهند. مثلاً بیت‌های ۱، ۰، ۱ مقدار ۵ دارند.
اما هر بردار منطقی الزاماً عدد نیست؛ می‌تواند سه سیگنال کنترل مستقل باشد.

\section{گیت‌های پایه و \lr{Truth Table}}
\subsection{تعریف عملگرها و ترتیب محاسبه}
\lr{Truth Table} خروجی را برای تمام ترکیب‌های مجاز ورودی می‌نویسد.
در این بخش ورودی‌ها فقط صفر یا یک هستند.
\lr{NOT} مقدار را معکوس می‌کند؛ آن را با خط بالای متغیر می‌نویسیم:
\(\overline{0}=1\) و \(\overline{1}=0\).
این خط علامت منفی عددی نیست.

عملگر \lr{AND} فقط وقتی یک است که همه ورودی‌ها یک باشند؛ نماد آن
\(A\cdot B\) یا \(AB\) است.
عملگر \lr{OR} وقتی یک است که دست‌کم یک ورودی یک باشد؛ نماد آن \(A+B\) است.
در اینجا \(+\) جمع معمولی نیست: \(1+1=1\) در جبر منطقی.
\lr{XOR} برای دو ورودی وقتی یک است که ورودی‌ها متفاوت باشند و با
\(A\oplus B\) نمایش داده می‌شود. \lr{NAND}، \lr{NOR} و \lr{XNOR}
به‌ترتیب مکمل \lr{AND}، \lr{OR} و \lr{XOR} هستند.
این تعریف‌ها زمینه خواندن گیت‌های بخش \secnum{۱.۴} کتاب را فراهم می‌کنند \cite{weste}.

جدول \ref{tab:foundation-gates} تمام حالت‌های دو ورودی را نشان می‌دهد.
برای جلوگیری از ابهام، در روابط ترکیبی از پرانتز استفاده می‌کنیم؛
در قرارداد متعارف، \lr{NOT} پیش از \lr{AND} و سپس \lr{OR} محاسبه می‌شود.
\begin{table}[ht]
\centering
\caption{جدول گیت‌های دوورودی؛ همه مدخل‌ها مقدار منطقی بدون واحد هستند.}
\label{tab:foundation-gates}
\begin{tabular}{cccccccc}
\toprule
\(A\)&\(B\)&\lr{AND}&\lr{OR}&\lr{XOR}&\lr{NAND}&\lr{NOR}&\lr{XNOR}\\
\midrule
۰&۰&۰&۰&۰&۱&۱&۱\\
۰&۱&۰&۱&۱&۱&۰&۰\\
۱&۰&۰&۱&۱&۱&۰&۰\\
۱&۱&۱&۱&۰&۰&۰&۱\\
\bottomrule
\end{tabular}
\end{table}
برای مثال، در سطر آخر، \lr{OR} یک ولی \lr{XOR} صفر است؛
پس «یا»ی \lr{OR} حالت هم‌زمان یک بودن هر دو ورودی را حذف نمی‌کند.

\subsection{مثال مرحله‌ای و ساخت جدول از عبارت}
برای عبارت \(Y=\overline{AB+C}\)، ابتدا \(AB\)، سپس \(AB+C\) و در پایان
مکمل کل عبارت را محاسبه می‌کنیم. با \(A=1\)، \(B=0\) و \(C=1\)،
داریم \(AB=0\)، سپس \(0+1=1\) و در نتیجه \(Y=0\).
خط مکمل فقط روی \(C\) نیست؛ محل و طول آن جزئی از تعریف تابع است.

برای ساخت جدول کامل، هر هشت ترکیب ورودی را می‌نویسیم و ستون میانی \(AB\)
را پیش از \(Y\) محاسبه می‌کنیم. راه فشرده‌تر برای فهم نتیجه این است:
اگر \(C=1\)، خروجی در همه حالت‌ها صفر است؛ اگر \(C=0\)،
داریم \(Y=\overline{AB}\). پس خروجی فقط در سه حالت
\((A,B,C)=(0,0,0),(0,1,0),(1,0,0)\) یک است.
این نتیجه با تعریف \lr{NAND} در جدول \ref{tab:foundation-gates} قابل کنترل است.

\subsection{از جدول به عبارت منطقی}
اگر خروجی تابعی فقط برای دو حالت \((A,B)=(0,1)\) و \((1,0)\) یک باشد،
برای حالت اول عبارت \(\overline{A}B\) و برای حالت دوم \(A\overline{B}\) را می‌نویسیم.
هر عبارت فقط همان حالت خود را تشخیص می‌دهد. سپس آن‌ها را با \lr{OR} ترکیب می‌کنیم:
\begin{equation}\label{eq:foundation-xor}
Y=\overline{A}B+A\overline{B}=A\oplus B.
\end{equation}
بنابراین جدول، عبارت و شبکه گیت‌ها سه نمایش مرتبط از یک تابع‌اند.
برای تعداد ورودی زیاد، روش‌های ساده‌سازی کارآمدتری وجود دارد؛
در این بخش هدف فهم نگاشت میان این نمایش‌هاست.

\section{جبر منطقی و اتصال \lr{Switch}ها}
\subsection{قانون‌های موردنیاز و کنترل آن‌ها}
برای متغیر دودویی \(A\)، روابط زیر از بررسی دو مقدار صفر و یک به دست می‌آیند:
\begin{align}
A+0&=A,& A\cdot1&=A,& A+1&=1,& A\cdot0&=0,\notag\\
A+A&=A,& AA&=A,& A+\overline{A}&=1,& A\overline{A}&=0.
\label{eq:foundation-identities}
\end{align}
مثلاً \(A+A=A\) نشان می‌دهد نباید جمع منطقی را مانند جمع عددی
به \(2A\) تبدیل کنیم.
دو قانون پرکاربرد دیگر، توزیع \(A(B+C)=AB+AC\) و جذب
\(A+AB=A\) هستند. برای جذب، اگر \(A=1\)، خروجی در هر صورت یک است؛
اگر \(A=0\)، هر دو جمله صفرند. پس دو طرف در همه حالت‌ها برابرند.

\lr{De Morgan's Laws} می‌گویند مکمل \lr{AND} برابر \lr{OR} مکمل‌ها و
مکمل \lr{OR} برابر \lr{AND} مکمل‌هاست \cite[بخش \secnum{۱.۴}]{weste}:
\begin{equation}\label{eq:foundation-demorgan}
\overline{AB}=\overline{A}+\overline{B},\qquad
\overline{A+B}=\overline{A}\,\overline{B}.
\end{equation}
برای رابطه اول، خروجی فقط وقتی صفر است که هر دو ورودی یک باشند؛
در غیر این صورت حداقل یک مکمل، یک است. این استدلال تمام چهار حالت را پوشش می‌دهد.
پس هنگام بردن مکمل داخل پرانتز، هم متغیرها مکمل می‌شوند و هم عملگر عوض می‌شود.

\textbf{مثال آموزشی.} با جذب داریم
\[
 F=A+\overline{A}B=(A+\overline{A})(A+B)=A+B.
\]
در این تبدیل از قانون توزیع دیگر جبر منطقی،
\(X+YZ=(X+Y)(X+Z)\)، استفاده شده است.
می‌توان بدون حفظ قانون نیز بررسی کرد: اگر \(A=1\)، هر دو طرف یک‌اند؛
اگر \(A=0\)، هر دو طرف برابر \(B\) هستند.

\subsection{چرا اتصال سری با \lr{AND} ارتباط دارد؟}
فعلاً \lr{Switch} ایدئالی را فرض می‌کنیم که با کنترل یک بسته می‌شود.
در دو \lr{Switch} سری، مسیر تنها وقتی بسته است که هر دو بسته باشند؛
پس شرط هدایت \(AB\) است. در دو \lr{Switch} موازی، بسته بودن یکی کافی است؛
پس شرط هدایت \(A+B\) است \cite[اسلاید ۲۲]{intro}.
این‌ها شرط هدایت مسیرند، نه الزاماً مقدار منطقی ولتاژ خروجی.
اگر شبکه مسیر خروجی به زمین را بسازد، برقرار شدن شرط هدایت می‌تواند خروجی را صفر کند.

در اتصال متعارف ریل‌به‌ریل، \lr{nMOS} با ورودی یک و \lr{pMOS} با ورودی صفر
هدایت می‌کند \cite[بخش \secnum{۱.۳}]{weste}\cite[اسلاید ۲۱]{intro}.
بنابراین برای \lr{pMOS} با کنترل \(A\)، شرط هدایت \(\overline{A}\) است.
این عبارت یک مدل مقدماتی برای اتصال مشخص است؛ شرط عمومی روشن شدن به اختلاف ولتاژ
\lr{Gate} و \lr{Source} و آستانه نیاز دارد.

\textbf{پل به تحلیل گیت.} دو \lr{nMOS} سری با کنترل \(A\) و \(B\)،
خروجی را وقتی \(AB=1\) به زمین وصل می‌کنند. برای داشتن گیت معتبر،
شبکه بالایی باید در بقیه حالت‌ها خروجی را به تغذیه برساند و در حالت اول قطع باشد.
با دو \lr{pMOS} موازی، شرط هدایت بالا
\(\overline{A}+\overline{B}=\overline{AB}\) است؛
پس تابع خروجی \(Y=\overline{AB}\)، یعنی \lr{NAND} خواهد بود.
طراحی شماتیک و تحلیل تفصیلی آن در بخش فصل اول دنبال می‌شود.

\section{مدارهای \lr{Combinational} و \lr{Sequential}}
\subsection{وابستگی به ورودی و وابستگی به گذشته}
در مدار \lr{Combinational}، پس از گذشت زمان لازم برای پایدار شدن،
خروجی از ورودی‌های همان زمان تعیین می‌شود. گیت‌های پایه و \lr{Multiplexer}
از نمونه‌های آن‌اند. وجود تاخیر فیزیکی، مدار را به‌خودی‌خود \lr{Sequential} نمی‌کند.

در مدار \lr{Sequential}، خروجی یا حالت بعدی علاوه بر ورودی، به حالت ذخیره‌شده
قبلی وابسته است. مثلاً شمارنده با ورودی کنترلی یکسان می‌تواند بسته به عدد قبلی،
عدد بعدی متفاوتی تولید کند. برای نمایش آن می‌توان نوشت:
\begin{equation}\label{eq:foundation-state}
q_{\mathrm{next}}=f(q_{\mathrm{current}},u).
\end{equation}
در اینجا \(q\) حالت، \(u\) ورودی و \(f\) قانون انتقال حالت است.
این رابطه تعریف ساختاری مسئله است، نه فرمول عددی ثابت برای همه مدارها.
مبانی مدارهای حافظه‌دار در بخش \secnum{۱.۴.۹} کتاب آمده‌اند \cite{weste}.

\subsection{نقش \lr{Clock} و تفاوت \lr{Latch} و \lr{Flip-Flop}}
\lr{Clock} سیگنالی زمان‌بندی‌شده است که در مدارهای هم‌زمان، زمان مجاز
به‌روزرسانی حالت را مشخص می‌کند. همه مدارهای حافظه‌دار الزاماً \lr{Clock} ندارند.
در مدل مقدماتی \lr{D Flip-Flop} لبه‌بالارونده، مقدار ورودی \(D\) در لبه
صفر به یک \lr{Clock} ثبت و تا لبه بعدی نگه داشته می‌شود؛ ورودی باید در
بازه زمانی لازم اطراف لبه پایدار باشد. این محدودیت‌ها در تحلیل دقیق‌تر با
\lr{Setup Time} و \lr{Hold Time} بیان می‌شوند.

\lr{D Latch} با فعال بودن \lr{Enable} شفاف است؛ خروجی می‌تواند تغییر ورودی را
دنبال کند و با غیرفعال شدن آن، مقدار را نگه دارد.
پس «حساس به سطح» بودن \lr{Latch} با «حساس به لبه» بودن \lr{Flip-Flop} تفاوت دارد
\cite[بخش \secnum{۱.۴.۹}]{weste}.
جزئیات ترانزیستوری آن‌ها بعد از معرفی گیت‌ها و \lr{Transmission Gate} آموزش داده می‌شود.

\textbf{مثال آموزشی.} اگر خروجی اولیه یک \lr{D Flip-Flop} صفر و مقدار
\(D\) در لبه اول یک باشد، خروجی پس از تاخیر انتشار یک می‌شود.
اگر \(D\) میان دو لبه به صفر تغییر کند، خروجی فوراً صفر نمی‌شود؛
با فرض رعایت محدودیت‌های زمانی، در لبه بعدی صفر ثبت خواهد شد.
مدل ثبت لبه‌ای را نباید با اتصال مستقیم \(Q=D\) اشتباه گرفت.

\section{تمرین‌های جمع‌بندی و پاسخ‌ها}
\subsection{ابتدا پاسخ خود را بنویسید}
\begin{enumerate}
\item اگر \(V_g=\qty{0.9}{\volt}\) و \(V_s=\qty{0.2}{\volt}\) باشد،
\(V_{gs}\) چقدر است؟ آیا فقط از این اطلاعات می‌توان روشن بودن ترانزیستور را نتیجه گرفت؟
\item در تقسیم ولتاژ، \(V_s=\qty{1.2}{\volt}\)، \(R_1=\qty{1}{\kilo\ohm}\)
و \(R_2=\qty{2}{\kilo\ohm}\) است. خروجی بی‌بار و جریان منبع را بیابید.
اگر بار \(\qty{2}{\kilo\ohm}\) به خروجی وصل شود، چه تغییری رخ می‌دهد؟
\item برای ظرفیت ثابت \(\qty{30}{\femto\farad}\)، تغییر ولتاژ از صفر به
\(\qty{0.6}{\volt}\) با جریان خالص ثابت \(\qty{9}{\micro\ampere}\) چقدر زمان می‌برد؟
\item تابع \(Y=\overline{A+B}\) در چه حالتی یک است؟ آیا با \(\overline{A}+\overline{B}\) برابر است؟
\item چرا مدار با تاخیر انتشار، الزاماً مدار حافظه‌دار نیست؟
\end{enumerate}

\subsection{پاسخ و کنترل نتیجه}
\begin{enumerate}
\item \(V_{gs}=V_g-V_s=\qty{0.7}{\volt}\).
برای نتیجه‌گیری درباره روشن بودن، نوع ترانزیستور، قرارداد ولتاژ و مقدار آستانه لازم است.
\item با رابطه \eqref{eq:foundation-divider}، خروجی بی‌بار
\(\qty{0.8}{\volt}\) و جریان منبع
\(1.2/3000=\qty{0.4}{\milli\ampere}\) است.
با بار، \(R_2\parallel R_L=\qty{1}{\kilo\ohm}\)،
خروجی \(\qty{0.6}{\volt}\) و جریان منبع \(\qty{0.6}{\milli\ampere}\) می‌شود.
کاهش مقاومت معادل بار باعث کاهش خروجی و افزایش جریان منبع شده است.
\item \(\Delta Q=C\Delta v=\qty{18}{\femto\coulomb}\).
بنابراین \(\Delta t=\Delta Q/I=\qty{2}{\nano\second}\).
واحد زمان و وابستگی نتیجه به فرض جریان ثابت باید کنترل شوند.
\item فقط برای \(A=B=0\) یک است.
طبق \lr{De Morgan} با \(\overline{A}\,\overline{B}\) برابر است؛
برای \((A,B)=(1,0)\)، عبارت اصلی صفر ولی
\(\overline{A}+\overline{B}\) یک می‌شود، پس دو عبارت اخیر برابر نیستند.
\item تاخیر انتشار فقط مدت رسیدن اثر ورودی به خروجی است.
حافظه‌دار بودن به وجود حالت ذخیره‌شده و وابستگی به گذشته نیاز دارد؛
یک \lr{Inverter} هم تاخیر دارد، اما در مدل منطقی پایدار \lr{Combinational} است.
\end{enumerate}

\subsection{آمادگی برای ورود به فصل اول}
اکنون باید بتوانید واحدها را تبدیل کنید، اختلاف ولتاژ و جهت جریان را مشخص کنید،
یک معادله گره ساده بنویسید، اثر خازن بر زمان تغییر را توضیح دهید و یک تابع منطقی
را از جدول یا اتصال \lr{Switch}ها استخراج کنید.
در گام بعد، این ابزارها برای فهم ساختار ترانزیستور، گیت‌های \lr{CMOS} و
مراحل طراحی فصل اول به کار خواهند رفت. هرجا ابزار تازه‌ای لازم باشد،
پیش از استفاده معرفی می‌شود.

برای حل یک مسئله تازه، ترتیب زیر را رعایت کنید:
\begin{enumerate}
\item گره‌ها، اتصال‌ها و مرجع ولتاژ را مشخص کنید؛ شباهت ظاهری شکل‌ها کافی نیست.
\item معلوم‌ها، مجهول‌ها و فرض مدل را بنویسید؛ مثلاً خروجی بی‌بار، ظرفیت ثابت
یا \lr{Switch} ایدئال. فرضی که داده نشده است باید صریح معرفی شود.
\item واحدها را یکسان کنید و سپس رابطه مناسب، مانند \lr{KCL}، قانون خازن
یا تابع منطقی، را بنویسید.
\item ابتدا رابطه را تا حد امکان نمادی حل کنید؛ سپس مقدارها را جای‌گذاری کنید.
\item جواب را در رابطه اولیه قرار دهید؛ واحد، علامت، بازه مجاز و رفتار حدی را کنترل کنید.
\end{enumerate}
برای نمونه، خروجی تقسیم ولتاژ با دو مقاومت مثبت و منبع مثبت باید میان صفر و منبع باشد؛
زمان شارژ با جریان خالص مثبت و تغییر ولتاژ مثبت نیز باید مثبت به دست آید.
چنین کنترل‌هایی جای محاسبه را نمی‌گیرند، اما خطاهای پنهان آن را آشکار می‌کنند.
